\par\Needspace{9\baselineskip}\bigskip
\pdfbookmark[3]{Exposition, [24]}{pallas-II-article24}
\begin{center}[24].\end{center}
\begingroup
\makeatletter
\def\@makefnmark{\hbox{\normalfont\@thefnmark}}
\makeatother
\noindent\hspace*{\parindent}Dans le cas d’un rapport rationel de \(n\text{ et }n'\), la valeur de \(n+\delta n\) contiendra une partie proportionnelle au tems ou une variation séculaire, et le moyen mouvement \(\int n\mathrm{d}t\) contiendra un terme proportionnel au quarré du tems. On peut s’en tenir pour les besoins actuels de l’astronomie, mais nous verrons que la considération des termes dépendants des puissances supérieures des masses rend une \textit{forme périodique} à cette équation séculaire en apparence. Cela est d’autant plus remarquable, parceque effectivement la théorie des mouvemens de Pallas va nous offrir ce cas\HistoricalNoteEight{Voir p. 557---559.] }.
\par\endgroup
